\section{Introduction}
\label{sec:intro}

Visual odometry and \ac{lidar}-based \ac{slam} are well-established, but both
degrade in dusty, foggy, featureless, or textureless environments.
Millimeter-wave (mmWave) \ac{fmcw} radar is robust to these
conditions: it penetrates smoke and dust, operates in the dark, and
measures Doppler velocity directly without feature matching.  These
properties make it attractive for aggressive flight in unstructured
environments.

Combining a mmWave radar with an \ac{imu} (\ac{rio}) for ego-motion estimation
has been explored with discrete-time \ac{ekf} and graph-optimization
formulations~\cite{doer2020ekf,kramer2020rve}.  However,
asynchronous radar-\ac{imu} measurements (radar at $\sim$10\,Hz, \ac{imu} at
$\sim$1\,kHz) motivate a continuous-time trajectory, a smooth function
evaluated at any sensor timestamp without interpolation or preintegration.

\textbf{Research question.}
This work asks: \emph{can a single-chip mmWave radar,
fused with an \ac{imu} and a magnetometer, provide usable 6-DoF 
state estimation, and if so, what does it take?}  The starting point
is what the sensors can and cannot do: per-return Doppler yields
direct ego-velocity, but at around a dozen returns per frame there
aren't enough features for position fixing.
The system is therefore a \emph{velocity and orientation} source whose
position estimate drifts with distance.
\autoref{sec:char} characterizes the radar measurements.
The rest of the report develops an estimator informed by this
characterization and evaluates it from gentle to aerobatic flight.

We contribute:
\begin{itemize}
  \item a per-return characterization of a single-chip radar from
        hover to $10$\,rad/s, with every residual measured against
        motion capture ground truth: an uncertainty law with a
        per-return and a per-frame component, so one configuration is
        used for every evaluated flight and platform.
  \item a continuous-time \ac{rio} pipeline based on that characterization:
        every radar return is an individually weighted Doppler factor in a B-spline estimator, with a consensus front
        end, sensor-only initialization, and a fixed-lag smoother with
        a marginalization prior, running in real time on a laptop CPU.
  \item an evaluation from gentle racing to aerobatic backflips, on
        held-out flights and on a public cross-platform dataset
        against published \ac{ekf} baselines, including a consistency
        check of the reported covariance.
\end{itemize}

\textbf{Data and code.}
The complete system is released for independent reproduction.  The datasets
and our Ceres solver with Python bindings are available at
\url{https://github.com/Dr1zz3l/spline-rio}.
